\clearpage\thispagestyle{empty}
\pdfbookmark[0]{Quellen zu Gauß an Schumacher (23. November 1847)}{schumacher18471123-sources}
\begin{center}\textbf{Quellen zu Gauß an Schumacher (23. November 1847)}\end{center}
\noindent Diese KI-Übertragung in heutiges Deutsch folgt dem deutschen Brief von Carl Friedrich Gauß an Heinrich Christian Schumacher vom 23.~November 1847 in den \href{https://gdz.sub.uni-goettingen.de/id/PPN236007467}{\textit{Werken}, Band~VI} (Göttingen, 1874), S.~194--196.
\par\smallskip\noindent Der \href{https://api.digitale-sammlungen.de/iiif/image/v2/bsb10049452_00009/full/1450,/0/default.jpg}{Erstdruck} erschien in den von H.~C.~Schumacher herausgegebenen \textit{Astronomischen Nachrichten}, Band~27, Nr.~625, Sp.~1--4. Das \href{https://api.digitale-sammlungen.de/iiif/image/v2/bsb10049452_00016/full/1450,/0/default.jpg}{Heft} ist auf den 29.~Januar 1848 in Altona datiert. Auf Gauß’ Unterschrift folgt dort die gesonderte Tafel \href{https://api.digitale-sammlungen.de/iiif/image/v2/bsb10049452_00010/full/1450,/0/default.jpg}{\textit{Zodiakus der Iris}} von Carl Wolfgang Benjamin Goldschmidt. Dessen \href{https://api.digitale-sammlungen.de/iiif/image/v2/bsb10049451_00128/full/1450,/0/default.jpg}{zweiter Satz von Bahnelementen der Iris} erschien in Band~26, Nr.~615, Sp.~239--240, im Heft vom 30.~Oktober 1847. Die Digitalisate der Zeitschrift stammen von der Bayerischen Staatsbibliothek.
\par\smallskip\noindent Übertragung und Satz wurden mit KI erstellt.
\clearpage
